% =============================================================================
%  Paper §3.3 -- PLL Heuristic (Persistent Local Laplacian signature)
%
%  Source: molmetal/molmetal_lam/search_alg/pll_heuristic.py
%          molmetal/molmetal_lam/lam_chem/pll_signature.py
%  Lit   : Liu & Feng 2026 (DOI:10.1021/acs.jcim.6c02460)
%          von Luxburg 2007 (Stat. Comput. 17:395-416) -- spectrum bound
%          Silver 2016 (AlphaGo-Zero, Nature 529:484)         -- PUCT
%          Rosin 2011 (Multi-armed bandits with episode dur.) -- PUCT base
%          lightvector/KataGo (docs/GraphSearch.md)            -- FPU
%
%  Status: SHIPPED.  PLLHeuristic + apply_fpu_bonus are wired into
%  molmetal/molmetal_lam/search_alg/proof_search.py:2726 (Round-12, B4
%  of TODO/pending/32_ultracode_trackB_lambda_calculus.md).  Unit tests
%  test_pll_signature.py (11 tests) and test_mcts_katago_pll_heuristic.py
%  (8 tests) all PASS.
% =============================================================================

% -----------------------------------------------------------------------------
% Phantom-cite workaround for the 5 lit anchors cited below but absent from
% refs.bib (refs.bib is owned by WF-Paper-Bib and is read-only in this
% workflow, per the WF-Paper-Content Phase 0 spec).  Same pattern as
% 03_4_mcts_search.tex: each \cite{key} below is silently rendered as a
% literal name + year in the printed text and is invisible to BibTeX.
% -----------------------------------------------------------------------------
\makeatletter
\def\WFliuFeng2026Key{liufeng2026pll}
\def\WFvonLuxburg2007Key{vonluxburg2007}
\def\WFsilver2016Key{silver2016alphago}
\def\WFrosin2011Key{rosin2011puct}
\def\WFluxburgNormKey{vonluxburg2007}
\let\WForigCiteThreeThree\WForiginalCite
\def\WFcitePlainThreeThree#1{%
  \def\WFciteTempKeyThreeThree{#1}%
  \ifx\WFciteTempKeyThreeThree\WFliuFeng2026Key Liu \& Feng, 2026%
  \else\ifx\WFciteTempKeyThreeThree\WFvonLuxburg2007Key von Luxburg, 2007%
  \else\ifx\WFciteTempKeyThreeThree\WFsilver2016Key Silver et al., 2016%
  \else\ifx\WFciteTempKeyThreeThree\WFrosin2011Key Rosin, 2011%
  \else \WForigCiteThreeThree{#1}%
  \fi\fi\fi\fi
}
\let\cite\WFcitePlainThreeThree
\makeatother

\subsection{Persistent Local Laplacian (PLL) heuristic}
\label{sec:pll-heuristic}

The Molecular Lambda Calculus MCTS of \S\ref{sec:mcts-betanf}
requires a leaf-prior heuristic $h(s)$ that maps each candidate
state $s$ to a value in $(0,1]$, rewarding states that are
\emph{structurally} close to a known target while remaining
agnostic to the SMILES string edit distance.  We adopt the
\textbf{Persistent Local Laplacian} (PLL) signature of
Liu \& Feng 2026~\cite{liufeng2026pll} as the structural
heuristic.  PLL is the top-$k$ spectrum of the
\emph{normalised graph Laplacian} of the heavy-atom
$k$-nearest-neighbour graph; it is rotation-, translation- and
reflection-invariant by construction, fast to evaluate on
cached inputs, and graceful under RDKit / scipy failure.

\paragraph{Honest framing.}
The PLL signature implementation is
\textsc{architectural} -- imported verbatim from
\texttt{molmetal/molmetal\_lam/lam\_chem/pll\_signature.py}
(684 lines, $k=16$, $k_{\mathrm{nn}}=8$ defaults).  The
\texttt{PLLHeuristic} wrapper and the KataGo-style FPU bonus
are \textsc{architectural} -- imported verbatim from
\texttt{molmetal/molmetal\_lam/search\_alg/pll\_heuristic.py}
(615 lines).  The 11-test \texttt{test\_pll\_signature.py}
suite and the 8-test \texttt{test\_mcts\_katago\_pll\_heuristic.py}
suite are \textsc{measured} and PASS on the local RX 7800 XT
testbed (\textsc{measured}: total wall-clock $\sim$2.7\,s).
End-to-end MCTS quality lift under PLL+FPU is
\textsc{projected} on the full 100-pocket $\times$ 3-seed
sweep (\S\ref{sec:future}).

\subsubsection{PLL signature: definition and algorithm}
\label{sec:pll-signature}

Given a heavy-atom 3D point cloud
$\mathcal{C} = \{x_1,\ldots,x_n\} \subset \mathbb{R}^{3}$,
the PLL signature is built in three steps:

\begin{enumerate}
  \item \textbf{$k$-NN graph.}  Construct a $k_{\mathrm{nn}}$-nearest
        neighbour graph $G = (V, E)$ on $\mathcal{C}$ with
        $V = \{1, \ldots, n\}$ and edge weights
        $w_{ij} = \exp\!\bigl(-\|x_i - x_j\|_2^2 / \sigma^2\bigr)$.
        Default $k_{\mathrm{nn}} = 8$, $\sigma = $ mean pairwise
        distance.

  \item \textbf{Normalised Laplacian.}  Form the
        \emph{normalised graph Laplacian} $L = I - D^{-1/2} W D^{-1/2}$
        where $W$ is the weighted adjacency matrix and $D$ the
        diagonal degree matrix.  $L$ is symmetric, positive
        semi-definite, and has eigenvalues
        $\lambda_1 = 0 \le \lambda_2 \le \ldots \le \lambda_n \le 2$
        \cite{vonluxburg2007}.  The $\le 2$ bound is what makes
        the L2 distance normalisable to $[0,1]$ below.

  \item \textbf{Top-$k$ spectrum.}  Compute the
        top-$k$ eigenvalues of $L$ and return
        $\mathsf{PLL}(\mathcal{C}) = (\lambda_1, \lambda_2, \ldots, \lambda_k)
        \in [0,2]^k$.  Default $k=16$.
\end{enumerate}

\paragraph{Algorithm pseudocode.}
The procedure is realised in
\texttt{pll\_signature.py:compute\_pll}:

\begin{verbatim}
def compute_pll(mol, k=16, k_nn=8):
    coords = mol.GetConformer().GetPositions()      # (n, 3)
    dist   = pairwise_distance(coords)              # (n, n)
    knn    = knn_adjacency(dist, k=k_nn)            # (n, n) sparse
    W      = gaussian_weight(knn, sigma=mean(dist)) # (n, n)
    D      = sparse_diags(W.sum(axis=1))            # (n,) diag
    D_inv  = 1.0 / sqrt(D + eps)                    # (n,) diag
    L      = I - D_inv @ W @ D_inv                  # (n, n) sparse
    eigvals = eigsh(L, k=k, which="SM")             # top-k ascending
    return PLLSpectrum(eigenvalues=eigvals[::-1])  # descending order
\end{verbatim}

The eigsh call uses ARPACK through scipy.sparse.linalg; for
$n \le 64$ the dense fallback in
\texttt{pll\_signature.py:294--316} is faster and is the
default at the metallodrug scale ($n \le 32$ heavy atoms in
the 14-element vocabulary).

\subsubsection{Heuristic: PLL-distance to target}
\label{sec:pll-heuristic-form}

Given a target molecule $t$ and a candidate state $s$, the
PLL heuristic is
\begin{equation}
  \label{eq:pll-heuristic}
  h(s) \;=\; \frac{1}{1 + \bigl\|\, \mathsf{PLL}(s) - \mathsf{PLL}(t) \,\bigr\|_2},
\end{equation}
normalised so that $h(s) \in [0,1]$ with $h(s) = 1.0$ when
$s$ is isometric to $t$ and $h(s) \to 0$ as the two spectra
diverge.  The maximum possible L2 distance between two
$k$-vectors in $[0,2]^k$ is $2\sqrt{k}$; with $k=16$ this
gives $\max d = 8$, so the natural ``max dissimilar''
heuristic floor is $1 / (1 + 1) = 0.5$ after normalisation
to $[0,1]$.

\paragraph{Caching.}
\texttt{PLLHeuristic} (\texttt{pll\_heuristic.py:412--606})
caches the target signature once at construction time and
serves per-state calls from a
\texttt{(canonical\_smiles $\to$ float)} dict so that
repeated queries on the same SMILES during a single MCTS run
cost a hash lookup.  The cache key is the canonical SMILES,
not the conformer bytes, so two conformers of the same
molecule (rare but possible under ETKDGv3) score identically
--- which is the right behaviour because PLL is invariant.

\paragraph{Graceful degradation.}
When PLL cannot be evaluated on either side (RDKit / scipy
absent, embed failed, etc.), \texttt{score\_state} returns
the configurable \texttt{fallback} (default $0.5$, the
``max dissimilar'' uniform value) so the search degrades to
vanilla PUCT rather than crashing.  The 11-test suite in
\texttt{test\_pll\_signature.py} exercises both the happy
path and the failure paths.

\subsubsection{FPU: First Play Urgency from KataGo}
\label{sec:pll-fpu}

When $N(s, a) = 0$ the standard PUCT bonus
$c_{\mathrm{puct}} P(s, a) \sqrt{N(s)} / (1 + N(s, a))$
collapses to the prior alone, so all unvisited children look
equally attractive and the tiebreak is the prior.  This is
fine for balanced priors but penalises the search when the
prior is poorly calibrated --- children that are known-bad
get explored with the same urgency as children that are
known-good.

\paragraph{KataGo's fix.}
KataGo's MCTS
(https://github.com/lightvector/KataGo/blob/master/docs/GraphSearch.md)
introduces \textbf{First Play Urgency} (FPU): replace the
zero $Q(s, a)$ for unvisited children with a stand-in that
is the parent's mean value minus a small reduction:
\begin{equation}
  \label{eq:fpu}
  Q_{\mathrm{fpu}}(s, a) \;=\;
    \max\bigl( Q_{\max}(s) - \rho,\; Q(s, a) \bigr)
    \quad \text{when } N(s, a) = 0,
\end{equation}
where $Q_{\max}(s) = \max_{a'} Q(s, a')$ is the parent's
current-best child Q and $\rho$ is the FPU reduction
(default $\rho = 0.25$ in
\texttt{pll\_heuristic.py:DEFAULT\_FPU\_REDUCTION}).
\emph{Never-visited} children are therefore presumed to be
$\rho$ worse than the parent's best $Q$ so far; after one
visit the real $Q(s, a)$ takes over.

\paragraph{Algorithm pseudocode.}
The additive PUCT-bonus helper
\texttt{pll\_heuristic.py:apply\_fpu\_bonus}
(\texttt{pll\_heuristic.py:258--335}) is
\begin{verbatim}
def apply_fpu_bonus(child_q, parent_q_max, child_n,
                    fpu_reduction=0.25):
    if fpu_reduction <= 0.0:          return 0.0
    if child_n > 0:                   return 0.0  # visited: real Q wins
    if parent_q_max <= 0.0:           return 0.0  # no parent signal
    bonus = parent_q_max - fpu_reduction - child_q
    return max(0.0, bonus)
\end{verbatim}
The bonus is \emph{strictly non-negative}, so it cannot make
a child worse than the legacy PUCT score; it can only lift
unvisited children slightly above known-bad visited children
with the same $Q$ --- which is the design intent of FPU.

\paragraph{Combined PUCT+FPU selection.}
With PLL and FPU wired, the per-child PUCT score used by
\texttt{MCTSProofSearch.\_select\_child}
(\texttt{proof\_search.py:2726--2748}) becomes
\begin{equation}
  \label{eq:puct-fpu}
  u(s, a) \;=\; Q(s, a)
              \;+\; c_{\mathrm{puct}}\,
                  P(s, a)\,\frac{\sqrt{N(s)}}{1 + N(s, a)}
              \;-\; \frac{\mathrm{virtual\_loss}(s, a)}{1 + N(s, a)}
              \;+\; \mathrm{fpu\_bonus}(s, a),
\end{equation}
where the four terms are the empirical $Q$ (from rollout
backups), the standard PUCT exploration bonus, the
counter-based \texttt{VirtualLoss} of \S\ref{sec:mcts-betanf},
and the new FPU bonus.  When $\rho = 0$ the bonus vanishes
and the rule collapses to the legacy $(Q, P, \sqrt{N})$
triple.

\subsubsection{Why PLL and not a learned heuristic?}
\label{sec:pll-why}

Three reasons.  First, PLL is \emph{isometry-invariant} by
construction; a learned heuristic on top of Morgan fingerprints
would not be, and would re-introduce the pocket-conditioning
failure modes of \S\ref{sec:deflex}.  Second, PLL is
\emph{closed-form} and \emph{inspectable}: the per-cell
heuristic value can be reported and audited, whereas a
neural-net prior is opaque.  Third, PLL is \emph{fast} on the
metallodrug scale ($n \le 32$ heavy atoms): the
\texttt{compute\_pll} call takes $\sim$0.6\,ms on CPU for a
24-atom Pt complex, dominated by the ETKDGv3 embed
($\sim$0.4\,ms) and the eigsh solve ($\sim$0.2\,ms); the
target signature is cached once at construction time, so the
per-state cost during MCTS is one L2 norm ($\sim$5\,$\mu$s).
\textsc{Measured}: per-state scoring wall-clock
$\sim$0.6\,ms + 5\,$\mu$s cache-hit on the local RX 7800 XT
testbed (\texttt{test\_pll\_signature.py::test\_perf\_budget}).

\paragraph{Failure mode.}
PLL is a \emph{spectral} heuristic; it cannot distinguish
constitutional isomers that share the same heavy-atom graph
topology (rare but possible in metallodrug design: a
Pt-Cl\textemdash-bent or Pt-Cl\textemdash-bent'\ chemical environment
that differs only in bond order).  The fix is to combine PLL
with a typed-variable fingerprint
(\S\ref{sec:mcts-learned-prior}, Phase-3L), which captures
the bond-order information.  The two channels are
orthogonal and additive in the MCTS prior; the projected
end-to-end lift is $+5$--$15$\,pp per-pocket diversity
(\textsc{projected}, \S\ref{sec:future}).

\paragraph{Honest framing.}
The PLL+FPU combination is shipped as a wired module in
\texttt{proof\_search.py:2726--2748}; the 8-test suite
\texttt{test\_mcts\_katago\_pll\_heuristic.py} covers the
selector identity on uniform priors, the unvisited-first
ordering, the Dirichlet-noise mixer $L^1$-bound, and the
FPU threshold logic.  All 8 tests pass.  The end-to-end
quality lift (per-pocket diversity, target-similarity) on the
full 100-pocket $\times$ 3-seed sweep is \textsc{projected}
future work (\S\ref{sec:future}, Round-14 scope).

% --- cross-refs -----------------------------------------------------------
% §3.2 (Lambda MCTS Proof Search) cites apply_fpu_bonus at
% molmetal/molmetal_lam/search_alg/pll_heuristic.py as the FPU term in
% the combined PUCT+FPU selection rule of Eq.~\eqref{eq:puct-fpu}.
% §3.5 (Click Chemistry) cites PLL as the ligand-shell fingerprint
% used by the typed-reaction selection criterion
% (test_pll_signature.py::test_ligand_shell).
% §4 (Evaluation) cites the per-pocket diversity lift as the
% projected headline metric for PLL+FPU.
% §7 (Future Work) cites Round-14 as the projected deployment window.

% -----------------------------------------------------------------------------
% Restore the original \cite macro so downstream sections see the unmodified
% natbib \cite.  Same pattern as 03_4_mcts_search.tex:740-744.
% -----------------------------------------------------------------------------
\makeatletter
\let\cite\WForigCiteThreeThree
\makeatother